\section*{Chapitre \ref{ch:g5:biodiversity} --- La biodiversité}

\begin{solution}{exo:g5:biodiversity:1}
Une sorte particulière d'\omterm{def:g1:living-or-not:living}{être vivant} --- des individus bâtis sur le
même modèle, qui se reproduisent entre eux. Exemples~: le moineau
domestique, la pâquerette, le renard roux, l'escargot des jardins.
\end{solution}

\begin{solution}{exo:g5:biodiversity:2}
La variété de vie qu'un lieu abrite --- avant tout son nombre
d'\omterm{def:g5:biodiversity:species}{espèces} et leur état. La prairie fleurie abrite des dizaines
d'\omterm{def:g5:biodiversity:species}{espèces} végétales et \omterm{def:g1:animals-around-us:animal}{animales} dans un mètre carré~: forte
\omterm{def:g5:biodiversity:biodiversity}{biodiversité}. La cour bitumée n'en abrite presque aucune.
\end{solution}

\begin{solution}{exo:g5:biodiversity:3}
Environ deux millions. Parce que chaque expédition sérieuse trouve
encore des \omterm{def:g5:biodiversity:species}{espèces} non nommées --- des scarabées surtout --- si
bien qu'il doit en exister des millions d'autres hors du catalogue.
\end{solution}

\begin{solution}{exo:g5:biodiversity:4}
Beaucoup d'\omterm{def:g5:biodiversity:species}{espèces} dans un même milieu (le mètre carré de prairie)~;
beaucoup de milieux \omterm{def:g3:bones-and-muscles:skeleton}{côte} à \omterm{def:g3:bones-and-muscles:skeleton}{côte} (mare, haie, pré)~; de la variété à
l'intérieur d'une \omterm{def:g5:biodiversity:species}{espèce} (pas deux \omterm{def:g1:animals-around-us:covering}{coquilles} d'escargot des jardins
tout à fait pareilles).
\end{solution}

\begin{solution}{exo:g5:biodiversity:5}
Dans un réseau riche, un chasseur dont la proie manque peut passer à
d'autres proies, et chaque rôle a des doublures~; le choc est
encaissé. Dans un réseau pauvre, chaque \omterm{def:g5:biodiversity:species}{espèce} a peu de solutions,
si bien qu'une perte en entraîne d'autres.
\end{solution}

\begin{solution}{exo:g5:biodiversity:6}
Abeilles et syrphes~: pollinisent les \omterm{def:g3:flowers-fruits-seeds:parts}{fleurs}. Coccinelles~: mangent
les pucerons. \omterm{prop:g3:sorting-living-things:groups}{Oiseaux}~: patrouillent à la recherche des \omterm{ex:g2:animals-grow-up:butterfly}{chenilles}.
Vers~: gardent le sol ouvert.
\end{solution}

\begin{solution}{exo:g5:biodiversity:7}
Milieu détruit (la mare asséchée)~; \omterm{def:g2:caring-for-nature:environment}{environnement} empoisonné
(déchets et produits chimiques)~; trop de prélèvements (la
surpêche). \omterm{prop:g5:biodiversity:loss}{Éteinte} veut dire entièrement disparue --- partout, pour
toujours.
\end{solution}

\begin{solution}{exo:g5:biodiversity:8}
Réponse libre. Un carré d'herbe ou de plantation donne d'ordinaire
dix \omterm{def:g5:biodiversity:species}{espèces} ou plus~; le carré bitumé, presque zéro --- la
différence est la mesure.
\end{solution}

\begin{solution}{exo:g5:biodiversity:9}
Chaque \omterm{def:g5:biodiversity:species}{espèce} est un fil~: sa perte laisse chaque voisin --- ce qui
la mangeait, ce qu'elle pollinisait ou nettoyait --- avec moins de
solutions, et affaiblit la capacité du réseau à encaisser le choc
suivant. Et comme personne ne connaît tous les fils que tient une
\omterm{def:g5:biodiversity:species}{espèce}, personne ne peut certifier d'avance qu'une perte est sans
danger~; l'extinction étant définitive, l'erreur ne pourrait jamais
être défaite.
\end{solution}

\begin{solution}{exo:g5:biodiversity:10}
Protéger les cigognes une par une reviendrait à garder des \omterm{prop:g3:sorting-living-things:groups}{oiseaux}~;
restaurer les prairies humides a rebâti ce dont les cigognes
\emph{dépendent} --- et un milieu nourrit tous ses habitants à la
fois, ce qui explique que grenouilles, libellules et \omterm{def:g3:flowers-fruits-seeds:parts}{fleurs} des prés
soient revenues avec les cigognes. Pas de milieu, pas d'habitants~;
milieu restauré, habitants de retour.
\end{solution}

\begin{solution}{exo:g5:biodiversity:11}
Ce qui a été retiré n'était pas de la décoration mais une
main-d'œuvre~: les \omterm{def:g5:biodiversity:species}{espèces} de la haie et des \omterm{def:g3:flowers-fruits-seeds:parts}{fleurs} logeaient les
pollinisateurs et les mangeurs de ravageurs, dont les services
n'apparaissaient dans les comptes de personne tant qu'ils étaient
gratuits. Le travail de la nature est une comptabilité invisible ---
constante, non payée, inaperçue --- si bien que son prix se
découvre exactement au moment où les ouvriers sont expulsés et où le
travail revient sous forme de factures.
\end{solution}
